% ButtonModule — ボタン入力モジュール
\subsection{ButtonModule}

\begin{apibox}{ButtonModule --- ボタン入力モジュール}
\textbf{定義ファイル:} \texttt{lib/ButtonModule/ButtonModule.h / .cpp}\\
\textbf{分類:} 入力モジュール（\texttt{updateInput()}のみオーバーライド）\\
\textbf{対象デバイス:} 汎用タクトスイッチ・プッシュボタン

GPIO入力をデバウンス処理し、押下状態とエッジ（立ち上がり・立ち下がり）を検出するモジュール。
プルアップ接続（activeLow）とプルダウン接続の両方に対応する。
\end{apibox}

\subsubsection{機能概要}

\begin{itemize}
    \item デバウンスフィルタ（\texttt{ModuleTimer}ベース、設定可能な安定時間）
    \item プルアップ（LOW=押下）/ プルダウン（HIGH=押下）の切り替え対応
    \item エッジ検出フラグ（\texttt{justPressed} / \texttt{justReleased}）
    \item 安定状態の保持（\texttt{pressed}）
\end{itemize}

% --- Config ---
\subsubsection{Config構造体}

\begin{funcblock}{ButtonConfig}
ボタンのハードウェア接続とデバウンスパラメータを定義する。

\begin{longtable}{l l l p{5.2cm}}
\toprule
\textbf{メンバ} & \textbf{型} & \textbf{制約・既定値} & \textbf{説明} \\
\midrule
\endhead
\texttt{pin}        & \texttt{uint8\_t}       & ---          & ボタン入力GPIOピン番号 \\
\texttt{activeLow}  & \texttt{bool}           & ---          & \texttt{true}: LOW=押下（プルアップ接続）、\texttt{false}: HIGH=押下（プルダウン接続） \\
\texttt{debounceMs} & \texttt{unsigned long}  & 例: 30       & デバウンス安定時間 [ms]。入力変化後、この時間が経過するまで状態を確定しない \\
\bottomrule
\end{longtable}
\end{funcblock}

% --- Data ---
\subsubsection{Data構造体}

\begin{funcblock}{ButtonData}
デバウンス処理後のボタン状態を保持する。\texttt{justPressed}/\texttt{justReleased}は
1ループ限りのエッジフラグであり、次の\texttt{updateInput()}呼び出しでクリアされる。

\begin{longtable}{l l l p{5.2cm}}
\toprule
\textbf{メンバ} & \textbf{型} & \textbf{初期値} & \textbf{説明} \\
\midrule
\endhead
\texttt{pressed}      & \texttt{bool} & \texttt{false} & デバウンス後の押下状態。押されている間\texttt{true} \\
\texttt{justPressed}  & \texttt{bool} & \texttt{false} & 今ループで押下された瞬間のみ\texttt{true}（立ち上がりエッジ） \\
\texttt{justReleased} & \texttt{bool} & \texttt{false} & 今ループで離された瞬間のみ\texttt{true}（立ち下がりエッジ） \\
\bottomrule
\end{longtable}
\end{funcblock}

% --- メソッド ---
\subsubsection{メソッド}

\begin{funcblock}{ButtonModule(const ButtonConfig\& config)}
\begin{tabularx}{\textwidth}{l X}
\toprule
\textbf{項目} & \textbf{内容} \\
\midrule
引数   & \texttt{config} --- \texttt{ButtonConfig}へのconst参照 \\
動作   & Configを\texttt{\_config}にコピー保持する \\
\bottomrule
\end{tabularx}
\end{funcblock}

\begin{funcblock}{bool init()}
\begin{tabularx}{\textwidth}{l X}
\toprule
\textbf{項目} & \textbf{内容} \\
\midrule
戻り値   & \texttt{bool} --- 常に\texttt{true} \\
動作     & \texttt{activeLow}に応じて\texttt{INPUT\_PULLUP}または\texttt{INPUT}でピンモードを設定。デバウンスタイマーを初期化 \\
\bottomrule
\end{tabularx}
\end{funcblock}

\begin{funcblock}{void updateInput(SystemData\& data)}
\begin{tabularx}{\textwidth}{l X}
\toprule
\textbf{項目} & \textbf{内容} \\
\midrule
書き込み先     & \texttt{data.button} \\
タイミング制御 & 毎ループ実行（タイマーはデバウンス判定に使用） \\
動作           & \texttt{digitalRead()}で生入力を取得し、前回値と比較。変化があればデバウンスタイマーをリセット。タイマー経過後に安定状態を確定し、\texttt{pressed}・\texttt{justPressed}・\texttt{justReleased}を更新 \\
\bottomrule
\end{tabularx}
\end{funcblock}

\begin{notebox}[エッジフラグのライフサイクル]
\texttt{justPressed}と\texttt{justReleased}は\texttt{updateInput()}の先頭で毎回\texttt{false}にリセットされる。
状態変化があったループでのみ\texttt{true}になるため、ロジックフェーズで1回だけ検出できる。
複数箇所で参照しても安全である（同一ループ内では値が変わらない）。
\end{notebox}

% --- 使用例 ---
\subsubsection{使用例}

\begin{lstlisting}[style=cppstyle, caption={ButtonModuleの使用例}]
// ProjectConfig.h
const ButtonConfig BUTTON_CONFIG = {
    .pin        = 0,     // GPIO0: BOOTボタン
    .activeLow  = true,  // LOW=押下（内蔵プルアップ使用）
    .debounceMs = 30,    // デバウンス30ms
};

// ロジックフェーズでの使用
void applyPattern(SystemData& data) {
    // ボタンが押された瞬間にLEDをトグル
    if (data.button.justPressed) {
        data.led.state = !data.led.state;
    }
}
\end{lstlisting}
